\clearpage
\fsection{Ej 7 Pág 153}
\fsubsection
[\textcolor{waveBlue2}{Ponte a prueba}]
{\textcolor{waveBlue2}{\boldit{Ponte a prueba}}}
{

	Encuentra los tres errores que aparecen en cada uno de los siguientes textos y redáctalos de manera correcta.

}



\begin{solution}
	\fsubsubsection{Texto 1}
	Una contraseña es una cadena de caracteres que nos permite \textcolor{samuraiRed}{\sout{identificarnos}} \textcolor{springGreen}{autenticarnos} para acceder a un dispositivo, programa o servicio. Para establecer una contraseña compleja debemos asegurarnos de que tenga, al menos, \textcolor{samuraiRed}{\sout{seis}} \textcolor{springGreen}{doce} caracteres que incluyan mayúsculas, minúsculas, números y caracteres especiales. Además, una contraseña debe ser aséptica, es decir, \textcolor{samuraiRed}{\sout{debe}} \textcolor{springGreen}{no debe} incluir información personal fácilmente identificable. Podemos utilizar un gestor de contraseñas que nos permita almacenar todas las contraseñas y proporcionárnoslas únicamente cuando accedemos al servicio con un correo electrónico.

	\vspace{2em}
	\hrulefill
	\vspace{2em}

	\fsubsubsection{Texto 2}
	El \textit{phishing} es un tipo de ataque que consiste en engañar a un usuario para que revele \textcolor{samuraiRed}{\sout{su correo electrónico}} \textcolor{springGreen}{sus credenciales o datos personales y bancarios}. Para ello, crean copias idénticas de páginas de acceso, correos electrónicos, etc., pero con enlaces maliciosos o formularios falsos, y las envían a la víctima, que comparte la información al confiar en la similitud con el contenido original. El método más común de \textit{phishing} es el \textcolor{samuraiRed}{\sout{\textit{whaling}}} \textcolor{springGreen}{\textit{spear phishing}}. Una de las medidas de prevención es la autenticación de dos factores, que implica proporcionar algo que se conoce (como una contraseña) y algo que se tiene (como un código generado por SMS o por \textcolor{samuraiRed}{\sout{un generador de contraseñas}} \textcolor{springGreen}{una aplicación de autenticación}).

\end{solution}
